Fix a declared feasible tuple and a feasible pair \((p,w)\), and define \(f\) and \(r_j\) as in the proposed statement.  The tail-sum definition and the zero-moment exactness equation give
\[
w_k=r_k,
\qquad
w_j=r_j-r_{j+1}\quad(1\le j<k),
\qquad
w_0=-r_1.
\tag{49}
\]
For \(1\le d\le\beta\), finite-sum integration and the exactness equations in Definition~\(\mathrm{def{:}unbiased\text{-}weight\text{-}class}\) give
\[
\begin{aligned}
\int_0^q d u^{d-1}f(u)\,du
&=\sum_{t=0}^kw_t\int_0^{p_t}d u^{d-1}\,du\\
&=\sum_{t=0}^kw_tp_t^d=1.
\end{aligned}
\tag{50}
\]
Let \(\phi_d(u)=d u^{d-1}\) in \(L^2[0,q]\).  Their Gram matrix is
\[
\langle\phi_d,\phi_e\rangle
=\int_0^qde\,u^{d+e-2}du
=R_\beta(q)_{d,e}.
\tag{51}
\]
The functions \(\phi_1,\ldots,\phi_\beta\) are linearly independent on every interval with \(q>0\).  Thus, for every nonzero \(b\in\mathbb R^\beta\),
\[
b^{\mathsf T}R_\beta(q)b
=\int_0^q\left\{\sum_{d=1}^\beta b_d\phi_d(u)\right\}^2du>0.
\]
This proves positive definiteness and justifies the inverse in Definition~\(\mathrm{def{:}degree\text{-}moment\text{-}barrier}\).

Put \(\lambda=R_\beta(q)^{-1}\mathbf1_\beta\) and
\[
g_\beta(u)=\sum_{e=1}^\beta\lambda_e\phi_e(u).
\]
Equation (51) gives \(\langle g_\beta,\phi_d\rangle=1\) for every \(d\), and (50) implies that \(f-g_\beta\) is orthogonal to the span of the \(\phi_d\)'s.  Hence
\[
\begin{aligned}
\int_0^q f(u)^2du
&=\int_0^q g_\beta(u)^2du
+\int_0^q\{f(u)-g_\beta(u)\}^2du\\
&\ge\lambda^{\mathsf T}R_\beta(q)\lambda
=\mathbf1_\beta^{\mathsf T}R_\beta(q)^{-1}\mathbf1_\beta.
\end{aligned}
\tag{52}
\]
This derives the Hilbert-space projection inequality from orthogonality.

The degree-one instance of (50) gives \(\int_0^qf=1\).  By Assumption~\(\mathrm{ass{:}common\text{-}uniform\text{-}rollout}\) and Lemma~\(\mathrm{lem{:}common\text{-}uniform\text{-}monomial\text{-}covariance}\), for a singleton \(S=\{i\}\),
\[
G(p,w)_{S,S}
=\operatorname{Var}\{f(U_i)\}
=\int_0^qf(u)^2du-1.
\tag{53}
\]
Combining (52)--(53) proves \(G(p,w)_{S,S}\ge\Psi_\beta(q)\).  The separately retained first-moment Cauchy--Schwarz inequality
\[
1=\left\{\int_0^qf(u)du\right\}^2
\le q\int_0^qf(u)^2du
\tag{54}
\]
gives \(G(p,w)_{S,S}\ge(1-q)/q\).  Applying the same Cauchy--Schwarz inequality to \(g_\beta\), for which \(\int_0^qg_\beta(u)du=1\), gives
\[
\mathbf1_\beta^{\mathsf T}R_\beta(q)^{-1}\mathbf1_\beta
=\int_0^qg_\beta(u)^2du\ge q^{-1},
\]
and therefore \(\Psi_\beta(q)\ge(1-q)/q\).  Since \(f=r_j\) on \((p_{j-1},p_j]\) and is zero on \((q,1]\), equations (50) with \(d=1\) and (53) also give
\[
\begin{aligned}
G(p,w)_{S,S}
&=\sum_{j=1}^k(p_j-p_{j-1})r_j^2-1\\
&=\sum_{j=1}^k(p_j-p_{j-1})(r_j-1)^2+(1-q)\\
&\ge\delta\sum_{j=1}^k(r_j-1)^2.
\end{aligned}
\tag{55}
\]

The singleton coordinate of \(\Omega^{1/2}G(p,w)\Omega^{1/2}\) is \(\omega_1G(p,w)_{S,S}\).  The coordinate Rayleigh quotient, Lemma~\(\mathrm{lem{:}subset\text{-}lattice\text{-}block\text{-}reduction}\), and Definitions~\(\mathrm{def{:}variance\text{-}certificate}\) and \(\mathrm{def{:}exact\text{-}saddle}\) imply
\[
V^*_{n,\beta,k,q,\delta}
\ge\frac{B^2\omega_1}{n}\Psi_\beta(q).
\tag{56}
\]

For \(\beta=1\), \(R_1(q)=(q)\), so
\[
\Psi_1(q)=q^{-1}-1=\frac{1-q}{q}.
\tag{57}
\]
For \(\beta\ge2\), the minimum-norm function under only the \(d=1\) constraint is \(q^{-1}\) on \([0,q]\), with squared norm \(q^{-1}\).  It satisfies the \(d=2\) constraint only if
\[
\int_0^q2u q^{-1}du=q=1.
\]
Thus, for \(q<1\), adding the second constraint strictly raises the squared projection norm.  Equations (52), (54), and the definition of \(\Psi_\beta\) give
\[
\Psi_\beta(q)>\frac{1-q}{q}.
\tag{58}
\]

For the small-budget expansion, write
\[
R_\beta(q)=D(q)\mathsf H_\beta D(q),
\qquad
D(q)_{d,d}=d q^{d-1/2},
\qquad
(\mathsf H_\beta)_{d,e}=\frac1{d+e-1}.
\tag{59}
\]
With \(b_d(q)=\{d q^{d-1/2}\}^{-1}\),
\[
\mathbf1_\beta^{\mathsf T}R_\beta(q)^{-1}\mathbf1_\beta
=b(q)^{\mathsf T}\mathsf H_\beta^{-1}b(q).
\tag{60}
\]
For fixed \(\beta\), the unique term with the largest negative power of \(q\) is the \((\beta,\beta)\) term.  Lemma~\(\mathrm{lem{:}hilbert\text{-}corner\text{-}inverse}\) therefore gives
\[
\Psi_\beta(q)
\sim\frac{(\mathsf H_\beta^{-1})_{\beta,\beta}}
{\beta^2q^{2\beta-1}}
=\frac{2\beta-1}{\beta^2}
\binom{2\beta-2}{\beta-1}^{\!2}q^{1-2\beta}.
\tag{61}
\]

Suppose \(\beta\ge2\) and \(q<1\).  Equality in (52) would require \(f=g_\beta\) almost everywhere.  The polynomial \(g_\beta\) cannot be constant: its first two moment constraints would otherwise force both its value to be \(q^{-1}\) and \(q=1\).  But \(f\) is constant on each interval \((p_{j-1},p_j)\), all of which have length at least \(\delta>0\).  A polynomial agreeing almost everywhere with a constant on a nonempty open interval is that constant polynomial everywhere.  Equality is impossible, and therefore
\[
G(p,w)_{S,S}>\Psi_\beta(q)
\tag{62}
\]
for every finite-stage feasible pair.  Lemma~\(\mathrm{lem{:}joint\text{-}design\text{-}attainment}\) supplies an attaining pair.  Applying the singleton-coordinate Rayleigh bound at it and using (62), together with \(B>0\) and \(\omega_1>0\), yields
\[
V^*_{n,\beta,k,q,\delta}
>\frac{B^2\omega_1}{n}\Psi_\beta(q).
\tag{63}
\]

If \(q=1\), fix arbitrary feasible nodes and use \(w_0=-1\), \(w_k=1\), with all intermediate weights zero.  These weights satisfy every exactness equation.  Every nonempty monomial is zero at \(p_0=0\) and one almost surely at \(p_k=1\), so its weighted readout is constant; hence \(G(p,w)=0\), \(V_{\mathrm{rob}}(p,w)=0\), and \(V^*_{n,\beta,k,1,\delta}=0\).  Conversely, a joint minimizer must have zero certificate.  The singleton coordinate bound gives \(G(p,w)_{S,S}=\operatorname{Var}\{f(U_i)\}=0\).  Since \(\mathbb Ef(U_i)=1\), \(f=1\) almost everywhere.  Positive interval lengths give \(r_j=1\) for every \(j\), and (49) forces the endpoint weights.  This proves the complete \(q=1\) characterization.

Finally let \(\beta=1\) and \(q<1\).  Distinct singletons have empty intersection, so Definition~\(\mathrm{def{:}exact\text{-}covariance\text{-}operator}\) gives \(G(p,w)_{S,T}=0\) for \(S\ne T\), while every diagonal entry is \(\operatorname{Var}\{f(U_i)\}\).  Since \(\Omega=\omega_1I\),
\[
V_{\mathrm{rob}}(p,w)
=\frac{B^2\omega_1}{n}\operatorname{Var}\{f(U_i)\}.
\tag{64}
\]
Equality in (54) occurs exactly when \(f=q^{-1}\) almost everywhere on \([0,q]\).  Positive interval lengths then give \(r_j=q^{-1}\) for every \(j\), and (49) gives the unique endpoint weights.  They are feasible for every node vector and attain variance \((1-q)/q\).  Equations (57) and (64) prove the final equality and show that every feasible node vector is jointly minimizing.